\documentclass[aspectratio=169,10pt,spanish]{beamer}

\input{../../../_preambulo_beamer.tex}
\input{../../../_comandos_beamer.tex}

\selectlanguage{spanish}

\title{MA1001B - Modelación Estadística para la Toma de Decisiones}
\subtitle{Lección 21: Distribución normal estándar y puntuaciones z}
\author{}
\institute{Tecnol\'ogico de Monterrey}
\date{}

\begin{document}

\begin{frame}[plain]
  \vspace{1.2cm}
  {\Large\bfseries MA1001B - Modelación Estadística para la Toma de Decisiones\par}
  \vspace{0.7cm}
  {\large Lección 21: Normal estándar y puntuaciones $z$\par}
  \vspace{0.7cm}
  {\color{TecRojo}\rule{\textwidth}{1.2pt}}
  \vspace{0.8cm}

  Facultad MA1001B\\[0.3cm]
  Periodo\\[0.3cm]
  Tecnol\'ogico de Monterrey
\end{frame}

\begin{frame}{La pregunta de estandarización}
  \begin{alertblock}{Pregunta guía}
    ¿Cómo puede una puntuación de auditoría de 60 en $N(50,8)$ leerse en la
    única regla normal estándar $Z\sim N(0,1)$?
  \end{alertblock}

  \vspace{0.6em}
  Al final de esta lección, usted puede:
  \begin{itemize}
    \item convertir una puntuación cruda en $z=(x-\mu)/\sigma$;
    \item leer $P(Z\le z)$ desde la curva estándar;
    \item transferir esa área de vuelta a $P(X\le x)$;
    \item explicar por qué el sesgo rompe una comparación de puntuación $z$.
  \end{itemize}

  \vfill
  \scriptsize Todas las puntuaciones de auditoría y tiempos de retraso usados hoy son sintéticos. No se usan datos reales de empresa.
\end{frame}

\begin{frame}{De una puntuación cruda a $z$}
  \small
  Modelo operativo: $X\sim N(50,8^2)$. La puntuación marcada es $x=60$:
  \[
    z=\frac{60-50}{8}=1.25.
  \]
  Una desviación estándar desde la media es siempre $z=\pm 1$:
  \[
    z(42)=-1,\qquad z(58)=1.
  \]

  \vspace{0.3em}
  \begin{exampleblock}{Reloj sintético de auditoría}
    $z$ localiza la puntuación en unidades de $\sigma$. No es en sí una
    probabilidad.
  \end{exampleblock}
\end{frame}

\begin{frame}[fragile]{Paso 1: Calcular $z$}
  \begin{columns}[T]
    \begin{column}{0.42\textwidth}
      \small
      \begin{lstlisting}
z = (x - mu) / sigma
z_60 = (60 - 50) / 8
      \end{lstlisting}

      \begin{block}{Salida verificada}
        $\mu=50$, $\sigma=8$\\
        $z(60)=1.250000$\\
        $z(\mu-\sigma)=-1.000000$\\
        $z(\mu+\sigma)=1.000000$
      \end{block}
    \end{column}
    \begin{column}{0.58\textwidth}
      \centering
      \includegraphics[width=\textwidth]{../figuras/normal_estandar_01_puntuaciones_z.png}
      \vspace{0.2em}
      \scriptsize Densidad operativa del Paso 1
    \end{column}
  \end{columns}
\end{frame}

\begin{frame}{La cola izquierda estándar}
  \small
  Después de la estandarización,
  \[
    P(X\le 60)=P(Z\le 1.25).
  \]
  \texttt{scipy.stats.norm.cdf(1.25)} devuelve $0.894350$. La cola derecha es
  \[
    P(Z>1.25)=1-0.894350=0.105650.
  \]
\end{frame}

\begin{frame}[fragile]{Paso 2: $P(Z\le 1.25)$}
  \begin{columns}[T]
    \begin{column}{0.40\textwidth}
      \small
      \begin{lstlisting}
z_model = norm(0, 1)
p_left = z_model.cdf(1.25)
      \end{lstlisting}

      \begin{block}{Salida verificada}
        $P(Z\le 1.25)=0.894350$\\
        $P(Z>1.25)=0.105650$\\
        $P(X\le 60)=0.894350$
      \end{block}
    \end{column}
    \begin{column}{0.60\textwidth}
      \centering
      \includegraphics[width=\textwidth]{../figuras/normal_estandar_02_probabilidades_estandar.png}
      \vspace{0.2em}
      \scriptsize Cola izquierda estándar del Paso 2
    \end{column}
  \end{columns}
\end{frame}

\begin{frame}{Paso 3: El sesgo rompe el percentil}
  \begin{columns}[T]
    \begin{column}{0.42\textwidth}
      \small
      Retrasos exponenciales: media $=\sigma=20$.
      El soporte empieza en $z=-1$.

      \vspace{0.4em}
      $P(Z\le -1)=0.158655$ en la campana.\\
      $P(X\le 0)=0.000000$ en el reloj de retraso.

      \vspace{0.4em}
      \textbf{Límite:} $z$ es comparable solo cuando las formas son similares.
      La Lección 22 usa cortes $z$ como límites de especificación.
    \end{column}
    \begin{column}{0.58\textwidth}
      \centering
      \includegraphics[width=\textwidth]{../figuras/normal_estandar_03_sesgo_rompe_z.png}
      \vspace{0.2em}
      \scriptsize Campana frente a reloj de retraso estandarizados del Paso 3
    \end{column}
  \end{columns}
\end{frame}

\begin{frame}{Verificación de conceptos}
  \begin{enumerate}
    \item Calcule $z$ para $x=60$ en $N(50,8)$.
    \item ¿Por qué $P(X\le 60)$ iguala $P(Z\le 1.25)$?
    \item ¿Qué probabilidad es $P(Z>1.25)$?
    \item ¿Por qué $z=-1$ no es el percentil 16 en el reloj de retraso?
  \end{enumerate}

  \vspace{0.6em}
  \begin{block}{Conexión con la Lección 22}
    La sesión puede continuar directamente con límites de especificación,
    rendimiento del proceso y un desplazamiento de la media.
  \end{block}

  \vfill
  \scriptsize Fuente: OpenStax, \textit{Introductory Business Statistics 2e},
  Capítulo 6, Sección 6.1. CC BY-NC-SA.
\end{frame}

\appendix



\end{document}
